%% fig_vickrey_bottleneck.tex
%% Vickrey ボトルネックモデル：出発時刻選択と最適動的課金
%% Usage: %% fig_vickrey_bottleneck.tex
%% Vickrey ボトルネックモデル：出発時刻選択と最適動的課金
%% Usage: %% fig_vickrey_bottleneck.tex
%% Vickrey ボトルネックモデル：出発時刻選択と最適動的課金
%% Usage: %% fig_vickrey_bottleneck.tex
%% Vickrey ボトルネックモデル：出発時刻選択と最適動的課金
%% Usage: \input{assets/assignment/fig_vickrey_bottleneck.tex}
%%
%% 図の構成
%%  左パネル（無課金均衡）：待ち時間カーブ，均衡費用=定数ライン
%%  右パネル（最適課金）  ：料金プロファイル（台形），待ち時間 = 0

\begin{tikzpicture}[
  every node/.style={font=\small},
  xscale=2.0,
  yscale=1.5
]

%% ============================================================
%% 共通パラメータ（正規化）
%%   ラッシュ開始 t=0, 希望到着時刻 t*=1, ラッシュ終了 t=2
%%   容量 s=1，需要 N=2（ラッシュ持続時間 = N/s = 2）
%%   α=1, β=0.5, γ=2
%% ============================================================

%% ============================================================
%% === 左パネル：無課金均衡（UE） ===
%% ============================================================

%% 軸
\draw[thick, ->] (-0.1, 0) -- (2.6, 0)
  node[below right] {出発時刻 $t$};
\draw[thick, ->] (0, -0.1) -- (0, 2.2)
  node[above left, align=center] {費用・待ち\\時間・課金};

%% 均衡費用（定数水平線）: C* = α·w + β·早着 or γ·遅着 が均等
%% 正規化: C*=1 とする
\draw[thick, blue!70!black, dashed]
  (0, 1.0) -- (2.4, 1.0)
  node[right] {$C^*$（均衡費用）};

%% 待ち時間 w(t)：出発時刻が早いほど待ち時間が長い（ラッシュ前半）
%%   前半 t∈[0,1]: w(t) 増加（行列が積み上がる）→ 簡略化：直線 (0,0.8)→(1,0)
%%   後半 t∈[1,2]: w(t) 減少
%% 三角形キュー: ピーク t*=1 で最大
\draw[thick, orange!80!black]
  (0, 0) -- (1, 0.5) -- (2, 0)
  node[right] {待ち時間 $w(t)$};

%% スケジュール遅延費用（早着/遅着）：鏡像のV字
\draw[thick, green!55!black]
  (0, 1.0) -- (1, 0.5) -- (2, 1.0)
  node[right, align=left, font=\footnotesize]
      {早着$\beta(\cdot)$\\[-3pt]+遅着$\gamma(\cdot)$};

%% 希望到着時刻 t* の垂直破線
\draw[gray, dashed, thin] (1, 0) -- (1, 1.1);
\node[below] at (1, 0) {$t^*$};

%% 軸ラベル
\node[below] at (0, 0) {$0$};
\node[below] at (2, 0) {$T$};

%% ラッシュ期間の括弧（\draw[decorate, decoration={brace}] の代わりに手動で描画）
\draw[thin] (0, -0.25) -- (0, -0.18);
\draw[thin] (2, -0.25) -- (2, -0.18);
\draw[thin] (0, -0.25) -- (2, -0.25);
  \node[below=3pt, font=\footnotesize] at (1, -0.25) {ラッシュ時間帯};

%% パネルタイトル
\node[font=\small\bfseries] at (1.2, 2.1) {（a）無課金均衡 --- キュー発生};

%% ============================================================
%% === 右パネル：最適課金 ===
%% ============================================================

\begin{scope}[xshift=3.2cm]

%% 軸
\draw[thick, ->] (-0.1, 0) -- (2.6, 0)
  node[below right] {出発時刻 $t$};
\draw[thick, ->] (0, -0.1) -- (0, 2.2)
  node[above left, align=center] {費用・待ち\\時間・課金};

%% 最適課金プロファイル τ(t)：台形型
%%   ラッシュ前半（t∈[0,1]）：増加   (0,0)→(1,0.5)
%%   ラッシュ後半（t∈[1,2]）：減少   (1,0.5)→(2,0)
\fill[red!15]
  (0, 0) -- (1, 0.5) -- (2, 0) -- cycle;
\draw[thick, red!70!black]
  (0, 0) -- (1, 0.5) -- (2, 0)
  node[right] {$\tau(t)$（最適課金）};

%% 最適課金後の費用（定数ライン = C* = 1）
\draw[thick, blue!70!black, dashed]
  (0, 1.0) -- (2.4, 1.0)
  node[right] {$C^*$（均衡費用）};

%% 待ち時間 = 0（キュー消滅）
\draw[thick, orange!80!black]
  (0, 0) -- (2, 0)
  node[right] {$w(t)=0$（キュー消滅）};

%% スケジュール遅延費用（V字，課金により C* に接する）
\draw[thick, green!55!black]
  (0, 1.0) -- (1, 0.5) -- (2, 1.0)
  node[right, align=left, font=\footnotesize]
      {スケジュール\\[-3pt]遅延費用};

%% 希望到着時刻 t*
\draw[gray, dashed, thin] (1, 0) -- (1, 1.1);
\node[below] at (1, 0) {$t^*$};

%% 軸ラベル
\node[below] at (0, 0) {$0$};
\node[below] at (2, 0) {$T$};

%% 課金収入領域ラベル
\node[font=\footnotesize, red!60!black, align=center]
  at (1.0, 0.18) {課金収入\\[-3pt]（三角形）};

%% パネルタイトル
\node[font=\small\bfseries] at (1.2, 2.1) {（b）最適動的課金 --- キュー消滅};

\end{scope}

\end{tikzpicture}

%%
%% 図の構成
%%  左パネル（無課金均衡）：待ち時間カーブ，均衡費用=定数ライン
%%  右パネル（最適課金）  ：料金プロファイル（台形），待ち時間 = 0

\begin{tikzpicture}[
  every node/.style={font=\small},
  xscale=2.0,
  yscale=1.5
]

%% ============================================================
%% 共通パラメータ（正規化）
%%   ラッシュ開始 t=0, 希望到着時刻 t*=1, ラッシュ終了 t=2
%%   容量 s=1，需要 N=2（ラッシュ持続時間 = N/s = 2）
%%   α=1, β=0.5, γ=2
%% ============================================================

%% ============================================================
%% === 左パネル：無課金均衡（UE） ===
%% ============================================================

%% 軸
\draw[thick, ->] (-0.1, 0) -- (2.6, 0)
  node[below right] {出発時刻 $t$};
\draw[thick, ->] (0, -0.1) -- (0, 2.2)
  node[above left, align=center] {費用・待ち\\時間・課金};

%% 均衡費用（定数水平線）: C* = α·w + β·早着 or γ·遅着 が均等
%% 正規化: C*=1 とする
\draw[thick, blue!70!black, dashed]
  (0, 1.0) -- (2.4, 1.0)
  node[right] {$C^*$（均衡費用）};

%% 待ち時間 w(t)：出発時刻が早いほど待ち時間が長い（ラッシュ前半）
%%   前半 t∈[0,1]: w(t) 増加（行列が積み上がる）→ 簡略化：直線 (0,0.8)→(1,0)
%%   後半 t∈[1,2]: w(t) 減少
%% 三角形キュー: ピーク t*=1 で最大
\draw[thick, orange!80!black]
  (0, 0) -- (1, 0.5) -- (2, 0)
  node[right] {待ち時間 $w(t)$};

%% スケジュール遅延費用（早着/遅着）：鏡像のV字
\draw[thick, green!55!black]
  (0, 1.0) -- (1, 0.5) -- (2, 1.0)
  node[right, align=left, font=\footnotesize]
      {早着$\beta(\cdot)$\\[-3pt]+遅着$\gamma(\cdot)$};

%% 希望到着時刻 t* の垂直破線
\draw[gray, dashed, thin] (1, 0) -- (1, 1.1);
\node[below] at (1, 0) {$t^*$};

%% 軸ラベル
\node[below] at (0, 0) {$0$};
\node[below] at (2, 0) {$T$};

%% ラッシュ期間の括弧（\draw[decorate, decoration={brace}] の代わりに手動で描画）
\draw[thin] (0, -0.25) -- (0, -0.18);
\draw[thin] (2, -0.25) -- (2, -0.18);
\draw[thin] (0, -0.25) -- (2, -0.25);
  \node[below=3pt, font=\footnotesize] at (1, -0.25) {ラッシュ時間帯};

%% パネルタイトル
\node[font=\small\bfseries] at (1.2, 2.1) {（a）無課金均衡 --- キュー発生};

%% ============================================================
%% === 右パネル：最適課金 ===
%% ============================================================

\begin{scope}[xshift=3.2cm]

%% 軸
\draw[thick, ->] (-0.1, 0) -- (2.6, 0)
  node[below right] {出発時刻 $t$};
\draw[thick, ->] (0, -0.1) -- (0, 2.2)
  node[above left, align=center] {費用・待ち\\時間・課金};

%% 最適課金プロファイル τ(t)：台形型
%%   ラッシュ前半（t∈[0,1]）：増加   (0,0)→(1,0.5)
%%   ラッシュ後半（t∈[1,2]）：減少   (1,0.5)→(2,0)
\fill[red!15]
  (0, 0) -- (1, 0.5) -- (2, 0) -- cycle;
\draw[thick, red!70!black]
  (0, 0) -- (1, 0.5) -- (2, 0)
  node[right] {$\tau(t)$（最適課金）};

%% 最適課金後の費用（定数ライン = C* = 1）
\draw[thick, blue!70!black, dashed]
  (0, 1.0) -- (2.4, 1.0)
  node[right] {$C^*$（均衡費用）};

%% 待ち時間 = 0（キュー消滅）
\draw[thick, orange!80!black]
  (0, 0) -- (2, 0)
  node[right] {$w(t)=0$（キュー消滅）};

%% スケジュール遅延費用（V字，課金により C* に接する）
\draw[thick, green!55!black]
  (0, 1.0) -- (1, 0.5) -- (2, 1.0)
  node[right, align=left, font=\footnotesize]
      {スケジュール\\[-3pt]遅延費用};

%% 希望到着時刻 t*
\draw[gray, dashed, thin] (1, 0) -- (1, 1.1);
\node[below] at (1, 0) {$t^*$};

%% 軸ラベル
\node[below] at (0, 0) {$0$};
\node[below] at (2, 0) {$T$};

%% 課金収入領域ラベル
\node[font=\footnotesize, red!60!black, align=center]
  at (1.0, 0.18) {課金収入\\[-3pt]（三角形）};

%% パネルタイトル
\node[font=\small\bfseries] at (1.2, 2.1) {（b）最適動的課金 --- キュー消滅};

\end{scope}

\end{tikzpicture}

%%
%% 図の構成
%%  左パネル（無課金均衡）：待ち時間カーブ，均衡費用=定数ライン
%%  右パネル（最適課金）  ：料金プロファイル（台形），待ち時間 = 0

\begin{tikzpicture}[
  every node/.style={font=\small},
  xscale=2.0,
  yscale=1.5
]

%% ============================================================
%% 共通パラメータ（正規化）
%%   ラッシュ開始 t=0, 希望到着時刻 t*=1, ラッシュ終了 t=2
%%   容量 s=1，需要 N=2（ラッシュ持続時間 = N/s = 2）
%%   α=1, β=0.5, γ=2
%% ============================================================

%% ============================================================
%% === 左パネル：無課金均衡（UE） ===
%% ============================================================

%% 軸
\draw[thick, ->] (-0.1, 0) -- (2.6, 0)
  node[below right] {出発時刻 $t$};
\draw[thick, ->] (0, -0.1) -- (0, 2.2)
  node[above left, align=center] {費用・待ち\\時間・課金};

%% 均衡費用（定数水平線）: C* = α·w + β·早着 or γ·遅着 が均等
%% 正規化: C*=1 とする
\draw[thick, blue!70!black, dashed]
  (0, 1.0) -- (2.4, 1.0)
  node[right] {$C^*$（均衡費用）};

%% 待ち時間 w(t)：出発時刻が早いほど待ち時間が長い（ラッシュ前半）
%%   前半 t∈[0,1]: w(t) 増加（行列が積み上がる）→ 簡略化：直線 (0,0.8)→(1,0)
%%   後半 t∈[1,2]: w(t) 減少
%% 三角形キュー: ピーク t*=1 で最大
\draw[thick, orange!80!black]
  (0, 0) -- (1, 0.5) -- (2, 0)
  node[right] {待ち時間 $w(t)$};

%% スケジュール遅延費用（早着/遅着）：鏡像のV字
\draw[thick, green!55!black]
  (0, 1.0) -- (1, 0.5) -- (2, 1.0)
  node[right, align=left, font=\footnotesize]
      {早着$\beta(\cdot)$\\[-3pt]+遅着$\gamma(\cdot)$};

%% 希望到着時刻 t* の垂直破線
\draw[gray, dashed, thin] (1, 0) -- (1, 1.1);
\node[below] at (1, 0) {$t^*$};

%% 軸ラベル
\node[below] at (0, 0) {$0$};
\node[below] at (2, 0) {$T$};

%% ラッシュ期間の括弧（\draw[decorate, decoration={brace}] の代わりに手動で描画）
\draw[thin] (0, -0.25) -- (0, -0.18);
\draw[thin] (2, -0.25) -- (2, -0.18);
\draw[thin] (0, -0.25) -- (2, -0.25);
  \node[below=3pt, font=\footnotesize] at (1, -0.25) {ラッシュ時間帯};

%% パネルタイトル
\node[font=\small\bfseries] at (1.2, 2.1) {（a）無課金均衡 --- キュー発生};

%% ============================================================
%% === 右パネル：最適課金 ===
%% ============================================================

\begin{scope}[xshift=3.2cm]

%% 軸
\draw[thick, ->] (-0.1, 0) -- (2.6, 0)
  node[below right] {出発時刻 $t$};
\draw[thick, ->] (0, -0.1) -- (0, 2.2)
  node[above left, align=center] {費用・待ち\\時間・課金};

%% 最適課金プロファイル τ(t)：台形型
%%   ラッシュ前半（t∈[0,1]）：増加   (0,0)→(1,0.5)
%%   ラッシュ後半（t∈[1,2]）：減少   (1,0.5)→(2,0)
\fill[red!15]
  (0, 0) -- (1, 0.5) -- (2, 0) -- cycle;
\draw[thick, red!70!black]
  (0, 0) -- (1, 0.5) -- (2, 0)
  node[right] {$\tau(t)$（最適課金）};

%% 最適課金後の費用（定数ライン = C* = 1）
\draw[thick, blue!70!black, dashed]
  (0, 1.0) -- (2.4, 1.0)
  node[right] {$C^*$（均衡費用）};

%% 待ち時間 = 0（キュー消滅）
\draw[thick, orange!80!black]
  (0, 0) -- (2, 0)
  node[right] {$w(t)=0$（キュー消滅）};

%% スケジュール遅延費用（V字，課金により C* に接する）
\draw[thick, green!55!black]
  (0, 1.0) -- (1, 0.5) -- (2, 1.0)
  node[right, align=left, font=\footnotesize]
      {スケジュール\\[-3pt]遅延費用};

%% 希望到着時刻 t*
\draw[gray, dashed, thin] (1, 0) -- (1, 1.1);
\node[below] at (1, 0) {$t^*$};

%% 軸ラベル
\node[below] at (0, 0) {$0$};
\node[below] at (2, 0) {$T$};

%% 課金収入領域ラベル
\node[font=\footnotesize, red!60!black, align=center]
  at (1.0, 0.18) {課金収入\\[-3pt]（三角形）};

%% パネルタイトル
\node[font=\small\bfseries] at (1.2, 2.1) {（b）最適動的課金 --- キュー消滅};

\end{scope}

\end{tikzpicture}

%%
%% 図の構成
%%  左パネル（無課金均衡）：待ち時間カーブ，均衡費用=定数ライン
%%  右パネル（最適課金）  ：料金プロファイル（台形），待ち時間 = 0

\begin{tikzpicture}[
  every node/.style={font=\small},
  xscale=2.0,
  yscale=1.5
]

%% ============================================================
%% 共通パラメータ（正規化）
%%   ラッシュ開始 t=0, 希望到着時刻 t*=1, ラッシュ終了 t=2
%%   容量 s=1，需要 N=2（ラッシュ持続時間 = N/s = 2）
%%   α=1, β=0.5, γ=2
%% ============================================================

%% ============================================================
%% === 左パネル：無課金均衡（UE） ===
%% ============================================================

%% 軸
\draw[thick, ->] (-0.1, 0) -- (2.6, 0)
  node[below right] {出発時刻 $t$};
\draw[thick, ->] (0, -0.1) -- (0, 2.2)
  node[above left, align=center] {費用・待ち\\時間・課金};

%% 均衡費用（定数水平線）: C* = α·w + β·早着 or γ·遅着 が均等
%% 正規化: C*=1 とする
\draw[thick, blue!70!black, dashed]
  (0, 1.0) -- (2.4, 1.0)
  node[right] {$C^*$（均衡費用）};

%% 待ち時間 w(t)：出発時刻が早いほど待ち時間が長い（ラッシュ前半）
%%   前半 t∈[0,1]: w(t) 増加（行列が積み上がる）→ 簡略化：直線 (0,0.8)→(1,0)
%%   後半 t∈[1,2]: w(t) 減少
%% 三角形キュー: ピーク t*=1 で最大
\draw[thick, orange!80!black]
  (0, 0) -- (1, 0.5) -- (2, 0)
  node[right] {待ち時間 $w(t)$};

%% スケジュール遅延費用（早着/遅着）：鏡像のV字
\draw[thick, green!55!black]
  (0, 1.0) -- (1, 0.5) -- (2, 1.0)
  node[right, align=left, font=\footnotesize]
      {早着$\beta(\cdot)$\\[-3pt]+遅着$\gamma(\cdot)$};

%% 希望到着時刻 t* の垂直破線
\draw[gray, dashed, thin] (1, 0) -- (1, 1.1);
\node[below] at (1, 0) {$t^*$};

%% 軸ラベル
\node[below] at (0, 0) {$0$};
\node[below] at (2, 0) {$T$};

%% ラッシュ期間の括弧（\draw[decorate, decoration={brace}] の代わりに手動で描画）
\draw[thin] (0, -0.25) -- (0, -0.18);
\draw[thin] (2, -0.25) -- (2, -0.18);
\draw[thin] (0, -0.25) -- (2, -0.25);
  \node[below=3pt, font=\footnotesize] at (1, -0.25) {ラッシュ時間帯};

%% パネルタイトル
\node[font=\small\bfseries] at (1.2, 2.1) {（a）無課金均衡 --- キュー発生};

%% ============================================================
%% === 右パネル：最適課金 ===
%% ============================================================

\begin{scope}[xshift=3.2cm]

%% 軸
\draw[thick, ->] (-0.1, 0) -- (2.6, 0)
  node[below right] {出発時刻 $t$};
\draw[thick, ->] (0, -0.1) -- (0, 2.2)
  node[above left, align=center] {費用・待ち\\時間・課金};

%% 最適課金プロファイル τ(t)：台形型
%%   ラッシュ前半（t∈[0,1]）：増加   (0,0)→(1,0.5)
%%   ラッシュ後半（t∈[1,2]）：減少   (1,0.5)→(2,0)
\fill[red!15]
  (0, 0) -- (1, 0.5) -- (2, 0) -- cycle;
\draw[thick, red!70!black]
  (0, 0) -- (1, 0.5) -- (2, 0)
  node[right] {$\tau(t)$（最適課金）};

%% 最適課金後の費用（定数ライン = C* = 1）
\draw[thick, blue!70!black, dashed]
  (0, 1.0) -- (2.4, 1.0)
  node[right] {$C^*$（均衡費用）};

%% 待ち時間 = 0（キュー消滅）
\draw[thick, orange!80!black]
  (0, 0) -- (2, 0)
  node[right] {$w(t)=0$（キュー消滅）};

%% スケジュール遅延費用（V字，課金により C* に接する）
\draw[thick, green!55!black]
  (0, 1.0) -- (1, 0.5) -- (2, 1.0)
  node[right, align=left, font=\footnotesize]
      {スケジュール\\[-3pt]遅延費用};

%% 希望到着時刻 t*
\draw[gray, dashed, thin] (1, 0) -- (1, 1.1);
\node[below] at (1, 0) {$t^*$};

%% 軸ラベル
\node[below] at (0, 0) {$0$};
\node[below] at (2, 0) {$T$};

%% 課金収入領域ラベル
\node[font=\footnotesize, red!60!black, align=center]
  at (1.0, 0.18) {課金収入\\[-3pt]（三角形）};

%% パネルタイトル
\node[font=\small\bfseries] at (1.2, 2.1) {（b）最適動的課金 --- キュー消滅};

\end{scope}

\end{tikzpicture}
